
\subsection{Data and sample}

This study uses data from the National Population Census of China conducted in
1982, 1990, 2000, and 2010.
Census data are widely regarded as the most comprehensive source of information
on the Chinese population, providing a large sample that covers all officially
recognized ethnic groups without sampling bias.

To identify married couples within individual-level Census data, we matched
spouses residing in the same household using five relationship categories to
the household head: head and spouse, parent and parent, parent-in-law and
parent-in-law, grandparent and grandparent, and children and children-in-law.
For households with multiple potential pairs within the same category, we used
year of first marriage (when available) as the matching criterion.
Because our matching method identified each marriage from both spouses' records,
we eliminated duplicates by focusing on female respondents.

Our final sample consists of 2,488,992 married women aged 25-34 who had
complete information for both their own and their spouse's ethnicity and
education.
We selected the 25-34 age range to minimize potential biases from
further educational advancement, early marriage, remarriage, bereavement, and
overlapping cohorts across census years.